\documentclass[11pt]{article}
\usepackage[T1]{fontenc}
\usepackage[utf8]{inputenc}
\usepackage{lmodern,amsmath,amssymb,microtype,xcolor}
\usepackage[margin=.9in]{geometry}
\usepackage[colorlinks=true,urlcolor=blue!60!black]{hyperref}
\setlength{\parindent}{0pt}
\setlength{\parskip}{.5em}
\setlength{\emergencystretch}{2em}
\author{Wentao Li\\\small Prepared with assistance from OpenAI Codex}
\date{November 22, 2026}
\title{Random forcing and oscillation of asymptotic density}
\hypersetup{pdfauthor={Wentao Li},pdftitle={Random forcing and oscillation of asymptotic density}}
\begin{document}
\maketitle
Brech, Brendle and Telles ask in Question 37 whether the least number of permutations needed to give every set of oscillating density a density is at least the covering number of the null ideal. Valderrama asks in Question 4.4 the parallel question for relative density on infinite subsets. Both questions have a negative consistency answer. The same model gives a counterexample to Conjecture 41 in the 2024 preprint and the 2025 journal article. \href{https://arxiv.org/abs/2410.21102}{Brech, Brendle and Telles, 2024, §5} · \href{https://doi.org/10.1112/jlms.70300}{Brech, Brendle and Telles, 2025, journal version} · \href{https://arxiv.org/abs/2506.21059}{Valderrama, 2025, §4}

\textbf{Definitions.} For infinite \(A\subseteq\omega\), \(d_A(X)\) denotes the limit of \(|X\cap A\cap n|/|A\cap n|\), and takes the symbol \(\mathrm{osc}\) if the limit fails to exist. Write \(d=d_\omega\). The cardinal \(\mathfrak r_{\rm osc}\) is the least size of a family of infinite sets \(A\) such that every infinite coinfinite \(X\) has \(d_A(X)\ne\mathrm{osc}\) for some member. The cardinal \(\mathfrak{dd}_{\{\mathrm{osc}\},\mathrm{all}}\) is defined using families of permutations which change the density of each oscillating set to a limit. The null covering number is \(\operatorname{cov}(\mathcal N)\).

\textbf{Theorem (relative to ZFC).} Starting from a model of GCH, add \(\aleph_2\) random reals by the product probability algebra. In the extension \[
\mathfrak{dd}_{\{\mathrm{osc}\},\mathrm{all}}=\mathfrak r_{\rm osc}=\aleph_1
<\operatorname{cov}(\mathcal N)=\aleph_2.
\] Moreover \(\mathfrak{dd}_{\{\mathrm{osc},1/2\},\mathrm{all}}\geq\aleph_2\), so replacing the input types \(\{\mathrm{osc}\}\) by \(\{\mathrm{osc},1/2\}\) changes the cardinal despite preserving the types outside \((0,1)\).

\textbf{Subsequence lemma.} For events \(E_n\) on a probability space there is a deterministic infinite coinfinite \(A=\{a_1<a_2<\cdots\}\) such that \(m^{-1}\sum_{j\le m}\mathbf1_{E_{a_j}}\) converges almost surely. The same choice works along every further deterministic subsequence, with a possibly different exceptional null set for each subsequence. This is the bounded Hilbert-space case of the classical Komlós subsequence principle. Here is an elementary proof. In the separable closed span of the indicators in \(L^2\), extract a weakly convergent subsequence with limit \(g\). It has \(0\le g\le1\) almost everywhere, by weak closedness of that convex set. Center by \(g\) and thin again to \(Y_j=\mathbf1_{E_{a_j}}-g\) with \(|\langle Y_i,Y_j\rangle|\le2^{-j}\) for all \(i<j\), skipping indices to leave infinite complement. For any further deterministic subsequence the mean of its first \(m\) terms has second moment at most \((m+2\sum_{j\ge1}j2^{-j})/m^2\). Chebyshev and Borel--Cantelli on the squares \(m=k^2\) give almost sure convergence to zero there. Boundedness and \((k+1)^2/k^2\to1\) fill the gaps.

\textbf{Forcing argument.} For a name \(\dot X\) for a subset of \(\omega\), choose measurable representatives \(E_n\) of the Boolean values of \(n\in\dot X\). Apply the lemma in the ground model. The forcing gives Boolean value one to existence of \(d_A(\dot X)\) for the resulting old \(A\), since a measurable convergence assertion holds almost everywhere. In particular the old infinite coinfinite subsets, of cardinality \(\aleph_1\), witness the upper bound for \(\mathfrak r_{\rm osc}\). The argument takes place separately for each name and needs no universal exceptional set.

For each such \(A\) use an old permutation mapping \(A\) in increasing order onto the complement of an infinite density-zero set, and map its complement bijectively onto the unused positions. If \(d_A(X)\) exists, the permuted set has the same ordinary density, since the unused positions have density zero. This proves the second upper bound directly.

No countable family can witness the relative-density invariant. Given \(A_0,A_1,\ldots\), form successive alternating included and excluded finite intervals whose right endpoints make each of \(A_0,\ldots,A_j\) contain more than \(j\) times its previous total within the new interval. Both relative-density extremes then occur for each \(A_i\). For the permutation invariant, given countably many \(\pi_i\) (include the identity), build a partial \(0\)--\(1\) assignment on larger finite domains. At stage \(j\) choose \(N\) so large that the old finite domain has size less than \(N/(j+1)\). Fill the previously unassigned points of \(\bigcup_{i\le j}\pi_i^{-1}[0,N)\) with the stage parity. The densities of all permuted sets oscillate between lower limit zero and upper limit one. Increase \(N\) with the stages and complete any remaining coordinates arbitrarily. These arguments give both lower bounds \(\aleph_1\).

Every real name for the product probability algebra depends modulo null sets on countably many coordinates. The algebra is ccc, so fewer than \(\aleph_2\) Borel null-set codes in the extension can be supported together on fewer than \(\aleph_2\) coordinates, by taking names for an enumeration and countable deciding antichains. A fresh random coordinate avoids all those null sets. Conversely there are at most \(\aleph_2\) reals by GCH and the countable-support calculation. This proves the asserted covering number.

Finally, in any model, fewer than \(\operatorname{cov}(\mathcal N)\) prescribed permutations cannot change the density of every set of density \(1/2\). For each fixed permutation the fair-coin strong law says the permuted random set still has density \(1/2\) almost surely. Include the identity and avoid the union of all the exceptional sets. This proves the last inequality and gives the claimed failure of Conjecture 41 with output types \(\mathrm{all}\). The use of one input family having empty intersection with \((0,1)\) is essential to this counterexample.

\textbf{Scope and priority.} The subsequence theorem, the covering calculation, and the permutation coding are classical or appear in the cited literature. The proposed contribution is the forcing application to the explicit density questions. The consistency separation gives a countermodel to each proposed universal inequality. Under CH their two sides are both \(\aleph_1\). Provisional significance \textbf{63/100}, conditional on correctness and novelty. It connects two independently formulated density questions to a standard obstruction in probability, within the theory of cardinal characteristics. Independent proof and priority review are pending. \href{https://doi.org/10.1007/BF02020976}{Komlós, 1967}
\end{document}
